\section{Approach}\label{sec:approach}

\subsection{Project Introduction To Problem Setting}
An IR system aims to retrieve relevant documents from a large corpus in response to a user's query. A fundamental approach to IR is the Vector Space Model (VSM), where both documents and queries are represented as vectors in a high-dimensional space defined by the vocabulary of the corpus. Similarity between a query and documents is typically computed using measures such as cosine similarity over TF-IDF weighted vectors.

In this project, we begin with a basic IR system built on the VSM framework using the Cranfield dataset. Retrieval is performed by ranking documents based on their similarity to the query.

While the VSM model provides a simple and computationally efficient framework, it relies on strong assumptions; most notably that words are independent and that meaning can be captured purely through term frequency statistics. These assumptions often lead to suboptimal retrieval performance, especially in the presence of word sense ambiguity, synonymy, and out-of-vocabulary (OOV) terms.

The goal of this project is to analyze the limitations of this baseline system and develop improved retrieval methods that better capture semantic relationships between words and documents. We evaluate system performance using standard IR metrics such as Precision, Recall, Mean Average Precision (MAP), and nDCG, and compare the effectiveness of enhanced models against the baseline VSM.

\subsection{Baseline}
The baseline system uses TF-IDF weighted unigram vectors and cosine similarity. It serves as the reference point against which all subsequent methods are compared.

\subsection{Some Other Approaches}

\subsubsection{Approach 1: Latent Semantic Analysis (LSA)}
To address the limitations of the Vector Space Model, specifically synonymy and polysemy, we implemented Latent Semantic Analysis (LSA). This approach moves away from exact keyword matching by uncovering the underlying semantic structure of the corpus using singular value decomposition.

\textbf{The Hypothesis:} Every document has underlying topics that are not explicitly stated but are reflected in word co-occurrences across the corpus. These co-occurrences are sufficient to capture synonymy and semantic context.

\textbf{Mathematical Formulation:}
\[
A = U \Sigma V^T, \qquad A \approx A_k = U_k \Sigma_k V_k^T
\]
where $U_k$ represents the term-concept mapping and $V_k^T$ represents the document-concept mapping in a reduced $k$-dimensional space.

\textbf{Implementation Details:}
\begin{enumerate}
    \item Log-TF weighting: $1 + \log(\text{freq})$.
    \item Smoothed IDF weighting before decomposition.
    \item Query projection into latent space:
    \[
    q_{latent} = q^T U_k \Sigma_k^{-1}
    \]
    \item Cosine similarity in the reduced latent space.
\end{enumerate}

\begin{table}[H]
\centering
\begin{tabular}{|l|c|}
\hline
\textbf{Process} & \textbf{Time (seconds)} \\ \hline
Index Building (SVD + Weighting) & 10.13 \\ \hline
Ranking (225 Queries) & 0.49 \\ \hline
\end{tabular}
\caption{LSA Execution Runtimes}
\end{table}

\begin{table}[H]
\centering
\small
\begin{tabular}{|c|ccc|ccc|}
\hline
\textbf{k} & \textbf{Precision} & \textbf{Recall} & \textbf{F-score} & \textbf{MAP} & \textbf{nDCG} & \textbf{MRR} \\ \hline
1  & 0.6222 & 0.1004 & 0.2840 & 0.6222 & 0.3174 & 0.6222 \\
2  & 0.5511 & 0.1741 & 0.3561 & 0.6778 & 0.3252 & 0.6778 \\
3  & 0.4993 & 0.2369 & 0.3795 & 0.6848 & 0.3365 & 0.6985 \\
4  & 0.4600 & 0.2858 & 0.3831 & 0.6849 & 0.3422 & 0.7074 \\
5  & 0.4213 & 0.3191 & 0.3726 & 0.6771 & 0.3467 & 0.7110 \\
6  & 0.3985 & 0.3581 & 0.3690 & 0.6662 & 0.3574 & 0.7169 \\
7  & 0.3810 & 0.3946 & 0.3650 & 0.6524 & 0.3662 & 0.7194 \\
8  & 0.3589 & 0.4225 & 0.3531 & 0.6424 & 0.3748 & 0.7222 \\
9  & 0.3407 & 0.4479 & 0.3425 & 0.6294 & 0.3840 & 0.7222 \\
10 & 0.3267 & 0.4752 & 0.3342 & 0.6186 & 0.3945 & 0.7235 \\ \hline
\end{tabular}
\caption{Evaluation Metrics for LSA Model on Cranfield Dataset}
\end{table}

\textbf{LSA-Synset hybrid approach:}
Another approach we explored in the LSA family was to disambiguate each word occurrence using its local context and WordNet, then build the LSA model on semantic concepts. For each token in a document or query, a lightweight Lesk-style algorithm selects the most appropriate WordNet synset. Once tokens are mapped to synsets, we build a concept-document matrix and proceed with standard LSA-based ranking.

\begin{table}[H]
\centering
\begin{tabular}{|c|c|c|c|c|c|c|}
\hline
K & Precision & Recall & F-score & MAP & nDCG & MRR \\ \hline
1  & 0.5333 & 0.0864 & 0.2431 & 0.5333 & 0.2844 & 0.5333 \\
2  & 0.4644 & 0.1476 & 0.2998 & 0.5844 & 0.2865 & 0.5844 \\
3  & 0.4281 & 0.1955 & 0.3208 & 0.5993 & 0.2900 & 0.6111 \\
4  & 0.3944 & 0.2371 & 0.3257 & 0.6041 & 0.2959 & 0.6233 \\
5  & 0.3618 & 0.2689 & 0.3181 & 0.6007 & 0.2999 & 0.6269 \\
6  & 0.3407 & 0.3003 & 0.3130 & 0.5919 & 0.3078 & 0.6313 \\
7  & 0.3213 & 0.3282 & 0.3061 & 0.5836 & 0.3146 & 0.6339 \\
8  & 0.3033 & 0.3516 & 0.2970 & 0.5827 & 0.3201 & 0.6367 \\
9  & 0.2859 & 0.3687 & 0.2859 & 0.5723 & 0.3233 & 0.6371 \\
10 & 0.2729 & 0.3895 & 0.2778 & 0.5644 & 0.3297 & 0.6389 \\ \hline
\end{tabular}
\caption{Evaluation Metrics for LSA-Synset Model on Cranfield Dataset}
\end{table}

\textbf{Analysis of Results and Hypothesis Evaluation:}
Our quantitative evaluation shows that LSA did not yield the expected performance boost over the baseline TF-IDF Vector Space Model. While it improved recall and some nDCG values, the overall gains were limited. Domain specificity, loss of fine-grained technical distinctions during dimensionality reduction, and the relatively small size of the Cranfield corpus likely limited the benefits of LSA. The LSA-synset hybrid further introduced semantic noise and WordNet mismatch with technical aeronautical terminology.

\subsubsection{Approach 2: N-gram Based Vector Space Model}
To capture local contextual information and meaningful phrases that are often lost in a unigram bag-of-words model, we implemented an n-gram retrieval system focused on bigrams ($n=2$).

\textbf{The Hypothesis:} Technical documents often contain multi-word expressions such as ``heat flow'' and ``boundary layer'' that carry more specific meaning together than as isolated words.

\textbf{Mathematical Formulation:}
\[
w_{g,D} = \text{tf}_{g,D} \times \log\left(\frac{N+1}{df_g+1}\right) + 1
\]
\[
\text{score}(D, Q) = \frac{\sum w_{g,D} \cdot w_{g,Q}}{\|V_D\| \|V_Q\|}
\]

\begin{table}[H]
\centering
\begin{tabular}{|l|c|}
\hline
\textbf{Process} & \textbf{Time (seconds)} \\ \hline
Index Building (Bigram generation + TF-IDF) & 0.38 \\ \hline
Ranking (225 Queries) & 0.50 \\ \hline
\end{tabular}
\caption{N-gram System Execution Runtimes}
\end{table}

\begin{table}[H]
\centering
\small
\begin{tabular}{|c|ccc|ccc|}
\hline
\textbf{k} & \textbf{Precision} & \textbf{Recall} & \textbf{F-score} & \textbf{MAP} & \textbf{nDCG} & \textbf{MRR} \\ \hline
1  & 0.5778 & 0.1002 & 0.2745 & 0.5778 & 0.3104 & 0.5778 \\
2  & 0.5067 & 0.1655 & 0.3323 & 0.6511 & 0.3156 & 0.6511 \\
3  & 0.4430 & 0.2101 & 0.3379 & 0.6567 & 0.3069 & 0.6659 \\
4  & 0.3756 & 0.2351 & 0.3146 & 0.6544 & 0.2951 & 0.6726 \\
5  & 0.3307 & 0.2551 & 0.2942 & 0.6511 & 0.2898 & 0.6761 \\
6  & 0.2978 & 0.2746 & 0.2772 & 0.6464 & 0.2889 & 0.6799 \\
7  & 0.2743 & 0.2917 & 0.2640 & 0.6299 & 0.2903 & 0.6811 \\
8  & 0.2572 & 0.3077 & 0.2537 & 0.6216 & 0.2936 & 0.6828 \\
9  & 0.2385 & 0.3167 & 0.2400 & 0.6143 & 0.2944 & 0.6833 \\
10 & 0.2249 & 0.3293 & 0.2305 & 0.6058 & 0.2975 & 0.6842 \\ \hline
\end{tabular}
\caption{Evaluation Metrics for N-gram Model ($n=2$)}
\end{table}

\textbf{Analysis:}
The bigram-based model performed lower than both the unigram VSM and the LSA approach. Tokenization and hyphenation issues, along with severe data sparsity in the bigram space, limited the effectiveness of this method.

\subsubsection{Approach 3: Sentence-wise Word Embedding Retrieval (Word2Vec)}
To move beyond exact term matching and latent topic modeling, we implemented a semantic retrieval system based on dense vector representations using Word2Vec embeddings.

\textbf{The Hypothesis:} By representing documents as sequences of sentence-level embeddings and using a soft-pooling mechanism, we can identify documents that contain specific semantic matches to query sentences regardless of exact vocabulary.

\textbf{Mathematical Formulation:}
\[
V_S = \frac{1}{|S|} \sum_{w \in S} v_w
\]
\[
\text{Score}(Q, D) = \frac{1}{\beta} \log \left( \frac{1}{n} \sum_{i=1}^{n} \exp(\beta \cdot \cos(V_Q, V_{S_i})) \right)
\]

\begin{table}[H]
\centering
\small
\begin{minipage}{0.48\textwidth}
\centering
\begin{tabular}{|c|ccc|c|}
\hline
\multicolumn{5}{|c|}{\textbf{Pretrained (Google News - 300d)}} \\ \hline
\textbf{k} & \textbf{Prec} & \textbf{Rec} & \textbf{MAP} & \textbf{nDCG} \\ \hline
1  & 0.4889 & 0.0810 & 0.4889 & 0.2796 \\
2  & 0.4222 & 0.1350 & 0.5556 & 0.2665 \\
3  & 0.3511 & 0.1640 & 0.5678 & 0.2447 \\
4  & 0.3211 & 0.1992 & 0.5790 & 0.2499 \\
5  & 0.2836 & 0.2179 & 0.5762 & 0.2454 \\
6  & 0.2600 & 0.2375 & 0.5678 & 0.2453 \\
7  & 0.2368 & 0.2493 & 0.5623 & 0.2445 \\
8  & 0.2228 & 0.2654 & 0.5457 & 0.2475 \\
9  & 0.2069 & 0.2749 & 0.5390 & 0.2488 \\
10 & 0.1973 & 0.2883 & 0.5301 & 0.2523 \\ \hline
\end{tabular}
\end{minipage}
\hfill
\begin{minipage}{0.48\textwidth}
\centering
\begin{tabular}{|c|ccc|c|}
\hline
\multicolumn{5}{|c|}{\textbf{Custom Trained (Cranfield - 100d)}} \\ \hline
\textbf{k} & \textbf{Prec} & \textbf{Rec} & \textbf{MAP} & \textbf{nDCG} \\ \hline
1  & 0.0578 & 0.0082 & 0.0578 & 0.0348 \\
2  & 0.0422 & 0.0116 & 0.0689 & 0.0280 \\
3  & 0.0459 & 0.0190 & 0.0774 & 0.0322 \\
4  & 0.0467 & 0.0243 & 0.0851 & 0.0332 \\
5  & 0.0444 & 0.0302 & 0.0914 & 0.0336 \\
6  & 0.0474 & 0.0396 & 0.0967 & 0.0379 \\
7  & 0.0489 & 0.0493 & 0.0986 & 0.0416 \\
8  & 0.0489 & 0.0575 & 0.1027 & 0.0438 \\
9  & 0.0479 & 0.0632 & 0.1050 & 0.0454 \\
10 & 0.0453 & 0.0663 & 0.1072 & 0.0460 \\ \hline
\end{tabular}
\end{minipage}
\caption{Rank-wise Comparison of Word2Vec Retrieval}
\end{table}

\textbf{Analysis:}
Both Word2Vec implementations underperformed relative to the baseline. The pretrained model suffered from domain mismatch, while the custom model was hurt badly by data scarcity. For a small technical corpus like Cranfield, exact term weighting remained more reliable than learned dense embeddings.

\subsubsection{Approach 4: Statistical Language Modeling (SLM) with Dirichlet Smoothing}
We also implemented a generative retrieval model based on Statistical Language Modeling (SLM) with Dirichlet smoothing.

\textbf{The Hypothesis:} A document is relevant to a query if the query is likely to have been generated by that document's language model.

\textbf{Mathematical Formulation:}
\[
\text{Score}(D, Q) = \sum_{w \in Q} \log P(w | M_d)
\]
\[
P(w | M_d) = \frac{c(w, D) + \mu \cdot P(w | C)}{|D| + \mu}
\]

\begin{table}[H]
\centering
\begin{tabular}{|l|c|}
\hline
\textbf{Process} & \textbf{Time (seconds)} \\ \hline
Index Building (Frequency Counts) & 0.0315 \\ \hline
Ranking (225 Queries) & 1.5194 \\ \hline
\end{tabular}
\caption{SLM Execution Runtimes}
\end{table}

\begin{table}[H]
\centering
\small
\begin{tabular}{|c|ccc|ccc|}
\hline
\textbf{k} & \textbf{Precision} & \textbf{Recall} & \textbf{F-score} & \textbf{MAP} & \textbf{nDCG} & \textbf{MRR} \\ \hline
1  & 0.6977 & 0.1179 & 0.3267 & 0.6977 & 0.3600 & 0.6977 \\
2  & 0.5577 & 0.1813 & 0.3637 & 0.7422 & 0.3398 & 0.7422 \\
3  & 0.5007 & 0.2437 & 0.3833 & 0.7366 & 0.3439 & 0.7555 \\
4  & 0.4411 & 0.2780 & 0.3684 & 0.7218 & 0.3391 & 0.7588 \\
5  & 0.4026 & 0.3117 & 0.3572 & 0.7134 & 0.3411 & 0.7677 \\
6  & 0.3622 & 0.3297 & 0.3352 & 0.7090 & 0.3387 & 0.7692 \\
7  & 0.3307 & 0.3469 & 0.3166 & 0.7008 & 0.3391 & 0.7705 \\
8  & 0.3100 & 0.3703 & 0.3053 & 0.6837 & 0.3422 & 0.7716 \\
9  & 0.2883 & 0.3840 & 0.2901 & 0.6749 & 0.3445 & 0.7721 \\
10 & 0.2724 & 0.4009 & 0.2791 & 0.6650 & 0.3500 & 0.7721 \\ \hline
\end{tabular}
\caption{Evaluation Metrics for SLM Model (Dirichlet Smoothing)}
\end{table}

\textbf{Analysis:}
SLM performed strongly because it preserved the benefits of exact term matching while addressing sparsity probabilistically through Dirichlet smoothing. It proved especially robust on this small technical corpus.

\subsubsection{Approach 7: BM25+}
To address the limitations of the baseline TF-IDF Vector Space Model, we implemented BM25+, an extension of the classical BM25 probabilistic ranking framework.

\textbf{Hypothesis:} BM25+ should outperform the baseline TF-IDF Vector Space Model because probabilistic term weighting and improved document length normalization better capture lexical relevance than cosine similarity-based retrieval.

\begin{table}[H]
\centering
\begin{tabular}{|c|c|c|c|c|c|c|}
\hline
k & Precision & Recall & F-score & MAP & nDCG & MRR \\
\hline
1 & 0.7022 & 0.1201 & 0.3304 & 0.7022 & 0.3770 & 0.7022 \\
2 & 0.5911 & 0.1980 & 0.3896 & 0.7467 & 0.3687 & 0.7467 \\
3 & 0.5185 & 0.2508 & 0.3964 & 0.7452 & 0.3646 & 0.7585 \\
4 & 0.4478 & 0.2828 & 0.3752 & 0.7384 & 0.3569 & 0.7630 \\
5 & 0.4080 & 0.3131 & 0.3614 & 0.7242 & 0.3570 & 0.7692 \\
6 & 0.3741 & 0.3405 & 0.3463 & 0.7153 & 0.3561 & 0.7736 \\
7 & 0.3416 & 0.3595 & 0.3275 & 0.7063 & 0.3554 & 0.7749 \\
8 & 0.3133 & 0.3739 & 0.3085 & 0.6989 & 0.3563 & 0.7749 \\
9 & 0.2938 & 0.3906 & 0.2954 & 0.6869 & 0.3597 & 0.7759 \\
10 & 0.2760 & 0.4045 & 0.2826 & 0.6763 & 0.3624 & 0.7763 \\
\hline
\end{tabular}
\caption{Evaluation Metrics for BM25+ Model on Cranfield Dataset}
\end{table}

\textbf{Analysis:}
BM25+ showed clear improvements over the baseline VSM on early precision, MAP, recall, and MRR. Better term-frequency saturation and document length normalization made it particularly effective for Cranfield's technical abstracts.

\subsubsection{Approach 8: BM25+ with Zero-Shot Neural Re-ranking}
To further improve retrieval performance beyond purely lexical matching, we implemented a hybrid retrieval pipeline combining BM25+ candidate retrieval with a transformer-based zero-shot neural reranker.

\begin{table}[H]
\centering
\begin{tabular}{|c|c|c|c|c|c|c|}
\hline
k & Precision & Recall & F-score & MAP & nDCG & MRR \\
\hline
1 & 0.6444 & 0.1082 & 0.2988 & 0.6444 & 0.3567 & 0.6444 \\
2 & 0.5733 & 0.1865 & 0.3736 & 0.7156 & 0.3567 & 0.7156 \\
3 & 0.5156 & 0.2479 & 0.3926 & 0.7237 & 0.3590 & 0.7363 \\
4 & 0.4567 & 0.2856 & 0.3806 & 0.7156 & 0.3531 & 0.7419 \\
5 & 0.4107 & 0.3193 & 0.3649 & 0.7062 & 0.3496 & 0.7463 \\
6 & 0.3785 & 0.3436 & 0.3503 & 0.6970 & 0.3510 & 0.7478 \\
7 & 0.3460 & 0.3627 & 0.3313 & 0.6880 & 0.3526 & 0.7484 \\
8 & 0.3239 & 0.3841 & 0.3185 & 0.6793 & 0.3561 & 0.7501 \\
9 & 0.2968 & 0.3933 & 0.2984 & 0.6722 & 0.3557 & 0.7501 \\
10 & 0.2760 & 0.4045 & 0.2826 & 0.6677 & 0.3571 & 0.7505 \\
\hline
\end{tabular}
\caption{Evaluation Metrics for BM25+ with Zero-Shot Neural Reranking}
\end{table}

\textbf{Analysis:}
BM25+ with zero-shot neural reranking did not yield the expected improvement over either the baseline TF-IDF Vector Space Model or standalone BM25+. Domain mismatch, semantic over-generalization, and added computational cost limited the practical benefit of reranking in this setting.

\subsubsection{Approach 9: Learning-to-Rank (LTR)}
To improve ranking quality, we implemented a Learning-to-Rank framework using a LambdaMART-based model.

\begin{table}[H]
\centering
\begin{tabular}{|c|c|c|c|c|c|c|}
\hline
k & Precision & Recall & F-score & MAP & nDCG & MRR \\
\hline
1 & 0.6618 & 0.1007 & 0.2911 & 0.6618 & 0.4301 & 0.6618 \\
2 & 0.5441 & 0.1708 & 0.3501 & 0.7059 & 0.3861 & 0.7059 \\
3 & 0.5049 & 0.2348 & 0.3804 & 0.7206 & 0.3777 & 0.7353 \\
4 & 0.4596 & 0.2776 & 0.3796 & 0.7157 & 0.3734 & 0.7426 \\
5 & 0.4176 & 0.3112 & 0.3680 & 0.6995 & 0.3711 & 0.7485 \\
6 & 0.3750 & 0.3251 & 0.3436 & 0.6965 & 0.3676 & 0.7534 \\
7 & 0.3466 & 0.3460 & 0.3290 & 0.6867 & 0.3699 & 0.7576 \\
8 & 0.3235 & 0.3675 & 0.3161 & 0.6871 & 0.3714 & 0.7595 \\
9 & 0.3023 & 0.3877 & 0.3025 & 0.6801 & 0.3759 & 0.7595 \\
10 & 0.2824 & 0.3974 & 0.2873 & 0.6753 & 0.3771 & 0.7595 \\
\hline
\end{tabular}
\caption{Evaluation Metrics for LTR Model on Cranfield Dataset}
\end{table}

\textbf{Analysis:}
The LTR framework did not strongly outperform the baseline in Precision@k, but it consistently achieved higher MAP, nDCG, and MRR across several ranks. This suggests it was somewhat better at ordering relevant documents higher within the ranked list.

\subsubsection{Approach 10: Learning-to-Rank with BM25-based Lexical Features}
To further improve retrieval effectiveness, we implemented an enhanced Learning-to-Rank system using BM25-derived lexical relevance features such as BM25 score, query term coverage, term frequency match strength, normalized document length, and query frequency.

\begin{table}[H]
\centering
\begin{tabular}{|c|c|c|c|c|c|c|}
\hline
k & Precision & Recall & F-score & MAP & nDCG & MRR \\
\hline
1 & 0.6029 & 0.0884 & 0.2596 & 0.6029 & 0.3419 & 0.6029 \\
2 & 0.5294 & 0.1655 & 0.3422 & 0.6618 & 0.3499 & 0.6618 \\
3 & 0.4461 & 0.2038 & 0.3383 & 0.6752 & 0.3203 & 0.6814 \\
4 & 0.3934 & 0.2369 & 0.3265 & 0.6626 & 0.3086 & 0.6814 \\
5 & 0.3559 & 0.2641 & 0.3153 & 0.6587 & 0.3065 & 0.6873 \\
6 & 0.3186 & 0.2830 & 0.2960 & 0.6481 & 0.3015 & 0.6873 \\
7 & 0.2983 & 0.3076 & 0.2868 & 0.6374 & 0.3056 & 0.6915 \\
8 & 0.2904 & 0.3437 & 0.2866 & 0.6052 & 0.3182 & 0.6915 \\
9 & 0.2778 & 0.3656 & 0.2804 & 0.6024 & 0.3261 & 0.6947 \\
10 & 0.2632 & 0.3823 & 0.2708 & 0.5925 & 0.3308 & 0.6947 \\
\hline
\end{tabular}
\caption{Evaluation Metrics for LTR + BM25 Feature Model}
\end{table}

\textbf{Analysis:}
Despite using more sophisticated lexical features, this enhanced LTR framework did not outperform the simpler TF-IDF baseline. Feature redundancy, small dataset size, strong baseline performance, and lack of deeper semantic features all contributed to the underwhelming gains.
